\section{Introduction}
\label{sec:introduction}

Urban Air Mobility (UAM) could become a new part of future transportation systems. To bring a large number of electric vertical take-off and landing (eVTOL) aircraft safely into low-altitude airspace, traffic-management frameworks such as the European U-space~\cite{eu_commission_2017,easa2020uspace} and the U.S. Unmanned Aircraft System (UAS) Traffic Management (UTM) concept~\cite{kopardekar2016utm} have been proposed. Both are built around services and trajectory-based operations~\cite{sesar_conops_2019,sesar_master_plan_2025}. The U-space framework itself does not define corridors. U-space services are provided in a certain region or volume, and UAS operators can plan their trajectories inside that volume. Still, some operators plan to use UAM corridors for the first UAM operations, and U-space services would then be provided for a larger volume that covers all corridors of a region. This setup is studied, for example, in the Single European Sky Air Traffic Management Research (SESAR) project EUREKA~\cite{sesar_eureka_project}. In this paper we look at such structured corridor operations within U-space; we do not treat U-space itself as a corridor concept.

For these operations, path planning has to be fast. In dense urban traffic, many trajectory requests must be handled while keeping separation and staying inside the corridors. Search-based methods such as A* have been studied extensively for UAS routing~\cite{sahawneh2016_3d_path_planning,a_star_2d3d_comparison_utm}, together with four-dimensional (4D) grid-based conflict management~\cite{acevedo_4d_grids_u_space}. Deconfliction in multi-dimensional trajectory spaces has also been studied for trajectory-based operations~\cite{NDMap,enea_porretta_4d_tbo_comparison}, and surveys give an overview of UAS path-planning methods~\cite{ahmad2023_uav_path_planning_survey}.

In an operational U-space setting, path planning is more than finding a route. The route must be turned into a 4D trajectory (position and time), checked against constraints, and prepared as a flight plan for approval. In this paper, we call this prepared plan the \emph{candidate operation intent}. It contains the origin and destination, the time window and the 4D trajectory, and an operator would submit it to a U-space Service Provider (USSP) for authorization. In the surrounding fleet-control architecture~\cite{nagel2026microservice}, this versioned operation-plan representation is termed a U-plan; the term is not taken from a regulatory definition. If all these functions sit in one large program, it is hard to extend them or to scale them separately. A microservice architecture splits them into separate software services, for example for the network topology, trajectory generation, deconfliction and plan submission~\cite{dragoni2017microservices,matlekovic2020uav}. However, there are only a few measurements of such service-based U-space planning systems. In particular, it is not well known how the planning time is split between the stages, how it depends on route length, and how a combined conflict-checking and re-planning loop behaves.

In this paper, we present a service-oriented 4D path-planning architecture for structured corridor operations within U-space and measure both the planning time per stage and the behavior of the system with conflict checking enabled. We use two existing tools of the German Aerospace Center (DLR). The Copter Flight Management System (CFMS) is an experimental flight-management and 4D trajectory generation system for unmanned multicopters and other vertical take-off and landing (VTOL) aircraft. It takes a waypoint sequence with flight phases and turns it into a dense, time-parameterized trajectory that is generated under the configured aircraft-performance assumptions~\cite{morscheck2021cfms}. NDMap (N-dimensional map) is a conflict-detection library for multi-dimensional objects. It stores trajectories in a combined space--time index and reports where they come too close to each other~\cite{NDMap}. In our architecture, route planning, waypoint synthesis, trajectory generation, deconfliction and the interface towards the USSP are separate service responsibilities. The Planning and Deconfliction Services only prepare a candidate trajectory; they do not authorize any flight.

The planning pipeline has four steps. First, A* finds the shortest corridor route on the vertiport graph. Second, waypoint synthesis adds departure, climb, cruise, descent and landing waypoints to this route. Third, CFMS turns the waypoints into a dense 4D trajectory. Fourth, if conflict checking is enabled, NDMap checks the new trajectory against the trajectories that were already accepted, and any conflicts are sent back to the Planning Service, which then tries another route (lateral re-routing). In this paper, one planning request means one operator request for a trajectory between one origin and one destination vertiport. The actual flight authorization by a USSP comes after this and is not part of the experiments.

We answer three research questions:
\begin{itemize}[leftmargin=3.6em,labelwidth=3em,align=left]
    \item[\textbf{RQ1}] How is the planning time split between A* route search, waypoint synthesis, and CFMS trajectory generation, and how does it change with route length on a fixed corridor network?
    \item[\textbf{RQ2}] When requests are handled one after another and the number of stored trajectories grows, how does the number of detected conflicts change when the same requests are packed into a shorter time window, and which of these conflicts can lateral re-routing solve?
    \item[\textbf{RQ3}] With conflict checking and re-routing enabled, where does the request time go---planning, the conflict-checking service call, or the remaining request overhead?
\end{itemize}

The contributions of this work are:
\begin{itemize}
    \item A service-oriented planning architecture for structured corridor operations within U-space, with route search, waypoint synthesis and trajectory generation inside the Planning Service and conflict detection in a separate Deconfliction Service.
    \item A waypoint-synthesis step that turns an A* corridor route into CFMS input with flight phases. We see this as an implementation step that makes the benchmark possible, not as a new algorithm.
    \item A breakdown of the planning time into A* search, waypoint synthesis and CFMS trajectory generation, and how it depends on route length (RQ1).
    \item A conflict-checking loop based on NDMap with lateral re-routing by removing graph edges and running A* again, tested with random and scripted encounter scenarios (RQ2), together with a breakdown of the total request time (RQ3).
\end{itemize}

The system described in this paper is an initial implementation. We evaluate it as a research prototype on a single machine, and the architecture, service boundaries and parameters are those of this first version. The operational system will be developed further and may differ from the setup shown here.

The rest of the paper is organized as follows. Section~\ref{sec:related} reviews related work; Section~\ref{sec:methodology} describes the system; Section~\ref{sec:experimental_setup} explains the evaluation setup; Section~\ref{sec:results} presents the results; Section~\ref{sec:discussion} discusses them and their limitations; and Section~\ref{sec:conclusion} concludes the paper.
